\documentclass[11pt, a4paper]{article}

% Gemeinsame Style-Definitionen (TEXINPUTS muss templates/ enthalten — siehe scripts/build_tex.py)
% bfw-style.tex — gemeinsame Style-Definitionen für alle Handouts/Aufgabenhefte
% Voraussetzung: Kompilation mit xelatex (wegen fontspec)

\usepackage[ngerman]{babel}
% Babel-ngerman macht " zu einem aktiven Shorthand (z.B. für "a -> ä). In unseren
% Texten verwenden wir " als normales Zeichen — daher Shorthand komplett ausschalten.
\shorthandoff{"}
\usepackage{geometry}
\geometry{a4paper, top=2.4cm, bottom=2.4cm, left=2.3cm, right=2.3cm,
          headheight=15pt, headsep=0.7cm, footskip=1.1cm}

\usepackage{fontspec}
% Fallback-Kette: Source Sans Pro -> Calibri -> Segoe UI -> Latin Modern Sans
% (Latin Modern ist immer mit MiKTeX vorhanden)
\IfFontExistsTF{Source Sans Pro}{%
  \setmainfont{Source Sans Pro}[Ligatures=TeX]%
}{\IfFontExistsTF{Calibri}{%
  \setmainfont{Calibri}[Ligatures=TeX]%
}{\IfFontExistsTF{Segoe UI}{%
  \setmainfont{Segoe UI}[Ligatures=TeX]%
}{%
  \setmainfont{Latin Modern Sans}[Ligatures=TeX]%
}}}
\IfFontExistsTF{Source Code Pro}{%
  \setmonofont{Source Code Pro}[Scale=0.92]%
}{\IfFontExistsTF{Consolas}{%
  \setmonofont{Consolas}[Scale=0.92]%
}{%
  \setmonofont{Latin Modern Mono}[Scale=0.92]%
}}

% Gesperrte Versalien für Labels/Titelbalken (Kapitälchen-Ersatz, funktioniert
% mit jeder OpenType-Schrift — Latin Modern Sans hat keine echten Kapitälchen)
\newcommand{\bfwlabelfont}{\footnotesize\bfseries\addfontfeatures{LetterSpace=8.0}}

\usepackage{xcolor}
% Farbpalette: hell, freundlich, modern
\definecolor{bfwPrimary}{HTML}{0EA5A4}    % Petrol/Türkis - Primär
\definecolor{bfwPrimaryDark}{HTML}{0F766E} % dunkler Petrol für Headings
\definecolor{bfwAccent}{HTML}{F59E0B}     % Amber - Akzent
\definecolor{bfwSuccess}{HTML}{10B981}    % Grün - Erfolg / Lösung
\definecolor{bfwWarn}{HTML}{F97316}       % Orange - Achtung
\definecolor{bfwInk}{HTML}{0F172A}        % fast Schwarz
\definecolor{bfwInkSoft}{HTML}{475569}    % gedämpftes Grau
\definecolor{bfwBgSoft}{HTML}{F8FAFC}     % off-white
\definecolor{bfwBgTeal}{HTML}{ECFEFF}     % helles Türkis
\definecolor{bfwBgAmber}{HTML}{FEF3C7}    % helles Gelb
\definecolor{bfwBgRose}{HTML}{FFE4E6}     % helles Rosé für Eselsbrücken
\definecolor{bfwTabHead}{HTML}{E4F3F2}    % zarte Kopfzeilen-Hinterlegung für Tabellen

\usepackage[colorlinks=true, linkcolor=bfwPrimaryDark, urlcolor=bfwPrimaryDark]{hyperref}
\usepackage{lastpage}
\usepackage{enumitem}
\setlist[itemize]{leftmargin=*, itemsep=2pt, topsep=2pt}
\setlist[enumerate]{leftmargin=*, itemsep=2pt, topsep=2pt}

\usepackage[most]{tcolorbox}
\usepackage{booktabs}
\usepackage{colortbl}
\usepackage{tikz}
\usetikzlibrary{decorations.pathmorphing, positioning, shapes.geometric, arrows.meta}

% ---------------------------------------------------------------------------
% Einheitliches Tabellen-Design
%   - etwas mehr Zeilenluft, dezente graue Linien (wirkt auch mit \hline ruhiger)
%   - Kopfzeile einer Tabelle bei Bedarf zart hinterlegen: \rowcolor{bfwTabHead}
% ---------------------------------------------------------------------------
\renewcommand{\arraystretch}{1.25}
\arrayrulecolor{bfwInkSoft!50}
\setlength{\heavyrulewidth}{1pt}
\setlength{\lightrulewidth}{0.5pt}

% ---------------------------------------------------------------------------
% Boxen — klare farbige Titelbalken mit gesperrten Versal-Labels
% (Titel bleibt per Option überschreibbar: \begin{merksatz}[title={...}])
% ---------------------------------------------------------------------------
\tcbset{bfwbox/.style={%
  enhanced, breakable,
  boxrule=0.5pt, arc=2.5pt,
  left=10pt, right=10pt, top=8pt, bottom=8pt,
  toptitle=4pt, bottomtitle=4pt,
  titlerule=0pt,
  fonttitle=\bfwlabelfont,
  coltitle=white
}}

% Merkkästchen — wichtige Definition / Kernpunkt
\newtcolorbox{merksatz}[1][]{%
  bfwbox,
  colback=bfwBgTeal,
  colframe=bfwPrimaryDark,
  colbacktitle=bfwPrimaryDark,
  title={MERKE},
  #1
}

% Eselsbrücke — Mnemoniks
\newtcolorbox{eselsbruecke}[1][]{%
  bfwbox,
  colback=bfwBgRose,
  colframe=bfwAccent!85!black,
  colbacktitle=bfwAccent!85!black,
  title={ESELSBRÜCKE},
  #1
}

% Beispiel-Box
\newtcolorbox{beispiel}[1][]{%
  bfwbox,
  colback=bfwBgSoft,
  colframe=bfwInkSoft,
  colbacktitle=bfwInkSoft,
  title={BEISPIEL},
  #1
}

% Lösungs-Box (für Aufgabenhefte)
\newtcolorbox{loesung}[1][]{%
  bfwbox,
  colback=bfwSuccess!7,
  colframe=bfwSuccess!65!black,
  colbacktitle=bfwSuccess!65!black,
  title={LÖSUNG},
  #1
}

% Achtung-Box — typische Fehler / Stolperfallen
\newtcolorbox{achtung}[1][]{%
  bfwbox,
  colback=bfwWarn!6,
  colframe=bfwWarn!85!black,
  colbacktitle=bfwWarn!85!black,
  title={ACHTUNG},
  #1
}

% Formel-Box — für zentrale Formeln (groß, ruhig, ohne Titelbalken)
\newtcolorbox{formelbox}[1][]{%
  enhanced,
  colback=bfwBgTeal,
  colframe=bfwPrimaryDark,
  boxrule=1pt, arc=3pt,
  left=10pt, right=10pt, top=6pt, bottom=6pt,
  #1
}

% Glossar-Eintrag
\newcommand{\glossareintrag}[2]{%
  \par\noindent\textbf{\color{bfwPrimaryDark}#1} \enspace #2\par\smallskip
}

% ---------------------------------------------------------------------------
% Abschnitts-Design: farbiges Nummern-Quadrat + feine Linie unter \section,
% dezent farbige \subsection/\subsubsection
% ---------------------------------------------------------------------------
\usepackage{titlesec}
\newcommand{\bfwSecBadge}[1]{%
  \begin{tikzpicture}[baseline={([yshift=-0.62em]n.center)}]
    \node[fill=bfwPrimaryDark, text=white, font=\normalsize\bfseries,
          minimum width=1.55em, minimum height=1.55em, inner sep=1pt,
          rounded corners=1.5pt] (n) {#1};
  \end{tikzpicture}}
\titleformat{\section}
  {\Large\bfseries\color{bfwPrimaryDark}}
  {\bfwSecBadge{\thesection}}{0.6em}{}
  [\vspace{-0.2em}{\color{bfwPrimary!60}\titlerule[0.8pt]}]
\titleformat{\subsection}
  {\large\bfseries\color{bfwPrimaryDark}}
  {\textcolor{bfwPrimary}{\thesubsection}}{0.5em}{}
\titleformat{\subsubsection}
  {\normalsize\bfseries\color{bfwInk}}
  {\textcolor{bfwPrimary}{\thesubsubsection}}{0.4em}{}
\titlespacing*{\section}{0pt}{1.6em}{1em}
\titlespacing*{\subsection}{0pt}{1.2em}{0.5em}

% TikZ-Stil "skizzenhaft" — etwas weicher als Rough.js, aber stabil
\tikzset{
  roughbox/.style = {
    draw=bfwPrimaryDark, thick, fill=bfwBgTeal,
    rounded corners=4pt, inner sep=6pt
  },
  roughaccent/.style = {
    draw=bfwAccent, thick, fill=bfwBgAmber,
    rounded corners=4pt, inner sep=6pt
  },
  roughline/.style = {
    draw=bfwInkSoft, thick, -{Stealth[length=2.5mm]}
  }
}

% Bullet-Symbole (bewusst kein FontAwesome — vermeidet Font-Installations-Probleme)
\newcommand{\faLightbulb}{$\bullet$}
\newcommand{\faBookmark}{$\bullet$}

% ---------------------------------------------------------------------------
% Kopf-/Fußzeile
%   Kopf links:  Wortlogo „Wirtschaftsfachwirt"
%   Kopf rechts: Dokument-Kurztext, gesetzt via \kopfzeile{...}
%   Fuß links:   © Can Siebert · 2026 — Fuß rechts: Seite X von Y
%   Titelseite (Seite 1) bleibt ohne Kopf-/Fußzeile (\handouttitel setzt
%   \thispagestyle{empty}).
% ---------------------------------------------------------------------------
\usepackage{fancyhdr}
\newcommand{\bfwKopfzeileText}{}
\newcommand{\kopfzeile}[1]{\renewcommand{\bfwKopfzeileText}{#1}}
\pagestyle{fancy}
\fancyhf{}
\fancyhead[L]{\small\bfseries\color{bfwPrimaryDark}Wirtschaftsfachwirt}
\fancyhead[R]{\small\color{bfwInkSoft}\bfwKopfzeileText}
\renewcommand{\headrulewidth}{0.6pt}
\renewcommand{\headrule}{{\color{bfwPrimary}\hrule width\headwidth height 0.6pt}}
\fancyfoot[L]{\small\color{bfwInkSoft}\copyright\ Can Siebert\,\textperiodcentered\,2026}
\fancyfoot[R]{\small\color{bfwInkSoft}Seite \thepage\ von \pageref*{LastPage}}
% Auch \thispagestyle{plain} (z.B. durch \maketitle) soll leer bleiben:
\fancypagestyle{plain}{\fancyhf{}\renewcommand{\headrulewidth}{0pt}\renewcommand{\headrule}{}}

% ---------------------------------------------------------------------------
% \handouttitel{Obertitel}{Untertitel}{Kurs-Label}{Termin-Zeile}
%   Einheitlicher professioneller Titelblock: Dachzeile, farbige Headline,
%   Untertitel, farbige Trennlinie und dezente Meta-Zeile (Kurs/Termin/Dozent).
%   Ersetzt \title/\maketitle. Direkt nach \begin{document} aufrufen.
% ---------------------------------------------------------------------------
\newcommand{\handouttitel}[4]{%
  \thispagestyle{empty}%
  \begingroup
  \setlength{\parindent}{0pt}%
  {\bfwlabelfont\color{bfwPrimary}WIRTSCHAFTSFACHWIRT (IHK)\par}
  \vspace{0.55em}
  {\huge\bfseries\color{bfwPrimaryDark}#1\par}
  \vspace{0.5em}
  {\large\color{bfwInkSoft}#2\par}
  \vspace{0.9em}
  {\color{bfwPrimary}\rule{\linewidth}{2pt}}\par
  \vspace{0.55em}
  {\small\color{bfwInkSoft}#3\ \,\textperiodcentered\,\ #4\ \,\textperiodcentered\,\ Dozent: Can Siebert\par}
  \endgroup
  \vspace{1.4em}%
}


\title{\vspace*{-1.5cm}\Huge\bfseries\color{bfwPrimaryDark}Lektion~3: Konjunktur, Wirtschaftspolitik \& Außenwirtschaft\\[0.4em]
  \large\color{bfwInkSoft}\textnormal{Wie EZB und Staat die Wirtschaft steuern --- und was die EU besonders macht}}
\author{\textsf{Wirtschaftsfachwirt --- 1.1 Volkswirtschaftliche Grundlagen \textperiodcentered{} \small\copyright\ Can Siebert\,\textperiodcentered\,2026}}
\date{}

\begin{document}
\maketitle
\thispagestyle{empty}

\vspace{1em}
\begin{merksatz}[title={\bfseries Worum geht es in dieser Lektion?}]
Die wirtschaftliche Entwicklung verläuft in Wellen (Konjunktur). Staat und Zentralbank
versuchen, diese Schwankungen zu glätten und die Ziele des Stabilitätsgesetzes zu erreichen.
In dieser Lektion lernen Sie die wichtigsten \textbf{wirtschaftspolitischen Maßnahmen}
kennen: die Geldpolitik der EZB, die Finanz-, Wachstums-, Tarif-, Arbeitsmarkt- und
Umweltpolitik --- sowie die beiden großen \textbf{Konzeptionen} (nachfrage- vs.
angebotsorientierte Wirtschaftspolitik). Im zweiten Teil geht es um die
\textbf{Außenwirtschaft}: Freihandel und Protektionismus sowie die Besonderheiten der EU
(Binnenmarkt, Wirtschafts- und Währungsunion). Damit schließen wir den Themenblock~1.1
des Rahmenplans ab --- alle Themen sind regelmäßig Gegenstand der IHK-Prüfung.
\end{merksatz}

\tableofcontents
\vspace{1em}
\hrule
\vspace{1em}

\section{Konjunktur: Zyklus und Indikatoren (Kurzüberblick)}

Jede volkswirtschaftliche Entwicklung ist Schwankungen unterworfen. Man unterscheidet
\textbf{Konjunkturschwankungen} (mehrjährige, wellenförmige Schwankungen der
gesamtwirtschaftlichen Aktivität um den langfristigen Wachstumstrend) von
\textbf{Saisonschwankungen} (regelmäßige, jahreszeitlich bedingte Schwankungen, z.\,B.
Bau im Winter, Tourismus im Sommer) und zufallsbedingten Schwankungen (z.\,B.
Naturkatastrophen, Krisen).

\subsection{Die vier Phasen des Konjunkturzyklus}

\begin{enumerate}
  \item \textbf{Aufschwung (Expansion):} steigende Nachfrage, Produktion und
    Investitionen, sinkende Arbeitslosigkeit, wachsender Optimismus.
  \item \textbf{Hochkonjunktur (Boom):} volle Kapazitätsauslastung, hohe Gewinne,
    Vollbeschäftigung, aber steigende Preise und Zinsen (Überhitzungsgefahr).
  \item \textbf{Abschwung (Rezession):} sinkende Nachfrage und Produktion,
    Auftragsrückgänge, zunehmende Arbeitslosigkeit, Pessimismus.
  \item \textbf{Tiefphase (Depression):} geringe Auslastung, hohe Arbeitslosigkeit,
    niedrige Preise/Zinsen --- Ausgangspunkt für einen neuen Aufschwung.
\end{enumerate}

\subsection{Konjunkturindikatoren}

\begin{itemize}
  \item \textbf{Frühindikatoren} (laufen der Entwicklung voraus): Auftragseingänge,
    Geschäftsklimaindex (ifo), Baugenehmigungen, Aktienkurse.
  \item \textbf{Präsenzindikatoren} (zeigen die aktuelle Lage): reales BIP,
    Kapazitätsauslastung, Industrieproduktion, Kurzarbeit, offene Stellen.
  \item \textbf{Spätindikatoren} (laufen hinterher): Arbeitslosenquote, Preisniveau
    (Inflationsrate), Insolvenzen, Lohnentwicklung.
\end{itemize}

\begin{eselsbruecke}
\textbf{F-P-S wie „Fahrplan der Konjunktur``:} \textbf{F}rüh kündigt an,
\textbf{P}räsenz zeigt an, \textbf{S}pät hinkt nach. Die Arbeitslosenquote ist ein
\emph{Spät}indikator --- Unternehmen entlassen erst, wenn der Abschwung schon läuft.
\end{eselsbruecke}

\section{Geldpolitik der Europäischen Zentralbank (EZB)}

Die Geldpolitik der EU liegt seit 1999 in den Händen der \textbf{Europäischen
Zentralbank (EZB)}. Sie steuert die Geldmengenentwicklung in der Eurozone und unterstützt
zudem die allgemeine Wirtschaftspolitik in der Gemeinschaft. \textbf{Vorrangiges Ziel}
ist nach dem Vertrag von Maastricht die \textbf{Preisniveaustabilität}: Der Anstieg des
\emph{harmonisierten Verbraucherpreisindex (HVPI)} soll mittelfristig \textbf{zwei
Prozent} betragen (seit der Strategieüberprüfung 2021 ein \emph{symmetrisches}
Inflationsziel; zuvor: „unter, aber nahe zwei Prozent``) --- damit sind sowohl Inflation
als auch Deflation ausgeschlossen. Die EZB
berücksichtigt bei ihrer Geldmengensteuerung sowohl die Preisentwicklung als auch die
Geldmengenentwicklung (\textbf{Zwei-Säulen-Strategie}).

\begin{merksatz}
Geschäftsbanken benötigen Zentralbankgeld, um Kredite vergeben zu können. Die Beschaffung
dieses Geldes heißt \textbf{Refinanzierung} --- wichtigste Refinanzierungsquelle ist die
Zentralbank. Genau hier setzen die geldpolitischen Instrumente an: Sie steuern
\emph{Liquidität} (Kreditmenge) und \emph{Zinsen} (Kreditkosten).
\end{merksatz}

\subsection{Instrument 1: Offenmarktgeschäfte}

Die Offenmarktgeschäfte stehen im Zentrum des Instrumentariums. Die EZB tritt als Käufer
oder Verkäufer von Wertpapieren am offenen Markt auf: \emph{Kauft} sie Wertpapiere, fließt
Zentralbankgeld in den Bankensektor (Geldschöpfung, Nachfrage wird belebt); \emph{verkauft}
sie, fließt Geld zurück (Geldvernichtung, Zinsen steigen, Nachfrage sinkt). Vier Kategorien:

\begin{itemize}
  \item \textbf{Hauptrefinanzierungsgeschäfte (Haupttender):} befristete Transaktionen im
    \emph{wöchentlichen} Rhythmus mit einer Laufzeit von \emph{einer Woche}
    (Pensionsgeschäfte bzw. Kredite gegen Verpfändung von Wertpapieren). Der hierfür
    festgelegte Zinssatz ist der \textbf{wichtigste Leitzins} im Euroraum.
  \item \textbf{Längerfristige Refinanzierungsgeschäfte (Basistender):} gleiche Technik,
    ursprünglich monatlich mit \emph{drei Monaten} Laufzeit (heute auch länger).
  \item \textbf{Feinsteuerungsoperationen:} gleichen kurzfristige Schwankungen der
    Bankenliquidität aus.
  \item \textbf{Strukturelle Operationen:} beeinflussen den Bedarf an Zentralbankgeld
    langfristig (z.\,B. Emission von Schuldverschreibungen).
\end{itemize}

Abgewickelt werden die Geschäfte in \textbf{Tenderverfahren} (standardisierte
Ausschreibungen): nach der Durchführung \emph{Standardtender} (24 Stunden, für Haupt- und
Basistender) oder \emph{Schnelltender} (wenige Stunden, für Feinsteuerung); nach der
Zuteilung \emph{Mengentender} (EZB gibt Volumen \emph{und} Zinssatz vor) oder
\emph{Zinstender} (Banken bieten Zinssätze, Zuteilung nach Höhe des Gebots).

\subsection{Instrument 2: Ständige Fazilitäten}

Die Fazilitäten stehen den Geschäftsbanken \emph{ständig} zur Verfügung (keine
Ausschreibung, Laufzeit: ein Geschäftstag):

\begin{itemize}
  \item \textbf{Spitzenrefinanzierungsfazilität} (Kreditfazilität): Banken können sich
    über Nacht unbegrenzt Zentralbankgeld leihen --- ihr Zinssatz bildet die
    \emph{Obergrenze} des Geldmarktzinses.
  \item \textbf{Einlagefazilität:} Banken legen überschüssige Liquidität bis zum nächsten
    Tag bei der EZB an --- ihr Zinssatz bildet die \emph{Untergrenze}.
\end{itemize}

Beide Sätze bilden zusammen den \textbf{Zinskorridor}, in dessen Mitte sich der
Hauptrefinanzierungssatz bewegt. Diese \textbf{drei Leitzinsen} (Einlagefazilität,
Hauptrefinanzierungssatz, Spitzenrefinanzierungsfazilität) müssen Sie für die Prüfung
nennen können.

\subsection{Instrument 3: Mindestreservepolitik}

Die EZB verpflichtet die Geschäftsbanken, bestimmte prozentuale Anteile der bei ihnen in
Sicht-, Termin- und Spareinlagen angelegten Gelder als \textbf{Pflichtguthaben} auf Konten
bei der Zentralbank zu unterhalten. Eine \emph{Erhöhung} des Mindestreservesatzes entzieht
den Banken Liquidität (weniger Kreditvergabe möglich), eine \emph{Senkung} wirkt umgekehrt.
Über Änderungen entscheidet der EZB-Rat.

\subsection{Expansive und restriktive Geldpolitik --- die Transmissionskette}

\begin{itemize}
  \item \textbf{Expansive Geldpolitik} („Politik des billigen Geldes``): Leitzinsen senken,
    Liquidität zuführen $\Rightarrow$ Kreditzinsen sinken $\Rightarrow$ Kreditnachfrage und
    Geldmenge steigen $\Rightarrow$ Investitionen und Konsum steigen $\Rightarrow$ BIP
    steigt, Arbeitslosigkeit sinkt --- Mittel gegen Rezession und Deflation.
  \item \textbf{Restriktive Geldpolitik}: Leitzinsen erhöhen, Liquidität verknappen
    (z.\,B. Mindestreserve erhöhen, Zuteilungsvolumen verringern) $\Rightarrow$ Zinsniveau
    steigt, Sparneigung nimmt zu $\Rightarrow$ kreditfinanzierte Investitionen und
    Konsumentenkredite gehen zurück $\Rightarrow$ gesamtwirtschaftliche Nachfrage sinkt
    $\Rightarrow$ Preisanstieg wird verringert, Konjunktur gedämpft --- Mittel gegen
    Inflation.
\end{itemize}

\begin{beispiel}
Die Inflationsrate liegt deutlich über zwei Prozent. Die EZB \emph{erhöht die Leitzinsen}:
Die Refinanzierung wird für die Geschäftsbanken teurer, sie geben die höheren Zinsen an
Unternehmen und Haushalte weiter. Kredite für Maschinen, Bauvorhaben und Konsum lohnen sich
weniger, die gesamtwirtschaftliche Nachfrage sinkt --- der Preisauftrieb lässt nach.
Kehrseite: Wachstum und Beschäftigung können leiden.
\end{beispiel}

\begin{eselsbruecke}
\textbf{„Zins runter --- Wirtschaft munter. Zins rauf --- Preisbremse drauf.``}
Expansiv = Gas geben (gegen Rezession), restriktiv = bremsen (gegen Inflation).
\end{eselsbruecke}

\subsection{Grenzen der Geldpolitik}

Die Instrumente wirken nur bedingt: Internationale Finanzströme und hohe Lohnabschlüsse
(über dem Produktivitätsfortschritt) können Inflation erzeugen; ob und wann Banken
Zinsänderungen weitergeben, ist unklar; bei pessimistischer Grundhaltung fördern selbst
niedrige Zinsen die Investitions- und Konsumbereitschaft nicht; und die EZB ist auf eine
einvernehmliche Fiskalpolitik der Regierungen angewiesen (erhöht der Staat in der
Inflationsbekämpfung gleichzeitig schuldenfinanziert seine Ausgaben, verpufft die
restriktive Wirkung).

\section{Finanzpolitik (Fiskalpolitik)}

Die Finanzpolitik umfasst alle Maßnahmen, die den \textbf{Staatshaushalt} (Einnahmen und
Ausgaben von Bund, Ländern, Gemeinden und Sozialversicherung) betreffen. Aufgaben:

\begin{itemize}
  \item \textbf{Beschaffung von Einnahmen} (v.\,a. Steuern, daneben Beiträge, Gebühren) zur
    Finanzierung öffentlicher Ausgaben; reichen sie nicht, steigt die Staatsverschuldung;
  \item \textbf{Umverteilung von Einkommen (Distribution)}, z.\,B. progressive
    Einkommensteuer, Transferleistungen, Sozialversicherung;
  \item \textbf{Stabilitätsaufgaben:} Steuerung der Konjunktur im Sinne des
    \textbf{Stabilitätsgesetzes (StabG, 1967)}.
\end{itemize}

\subsection{Antizyklische Fiskalpolitik nach Keynes}

Nach der Lehre von \textbf{John Maynard Keynes} sollen staatliche Einnahmen- und
Ausgabenpolitik dem Konjunkturverlauf \emph{entgegenwirken}:

\begin{center}
\begin{tabular}{p{3.2cm} p{5.2cm} p{5.2cm}}
\toprule
 & \textbf{Boom (restriktiv)} & \textbf{Rezession (expansiv)}\\
\midrule
Ausgaben & senken (Investitionen, Subventionen, Transfers zurückfahren) & erhöhen (öffentliche Aufträge, Konjunkturprogramme)\\
Einnahmen & erhöhen (Steuern) & senken (Steuererleichterungen)\\
Budgetwirkung & Überschüsse bilden (Rücklagen) & Defizite zulassen (\emph{Deficit Spending})\\
\bottomrule
\end{tabular}
\end{center}

Nimmt der Staat in der Rezession bewusst neue Schulden auf, um zusätzliche Nachfrage zu
erzeugen, spricht man von \textbf{Deficit Spending}.

\begin{beispiel}
In der Finanz- und Wirtschaftskrise wurden die \emph{Konjunkturpakete I und II} (2008/2009)
verabschiedet: Investitionen in Infrastruktur und öffentliche Gebäude,
Abschreibungserleichterungen (wirkten wie eine Steuererleichterung) und die
\emph{Umweltprämie (Abwrackprämie)}, die die Nachfrage nach Kraftfahrzeugen fördern sollte.
\end{beispiel}

\subsection{Automatische Stabilisatoren}

Neben gezielten Maßnahmen wirken \textbf{automatische (eingebaute) Stabilisatoren} ohne
politische Entscheidung antizyklisch: Die \emph{progressive Einkommensteuer} entzieht im
Boom überproportional Kaufkraft und entlastet in der Rezession; \emph{Arbeitslosengeld und
Sozialleistungen} stützen in der Rezession die Nachfrage, während im Aufschwung die
Beiträge die Kaufkraft dämpfen. Der Haushalt „atmet`` also mit der Konjunktur.

\subsection{Staatsverschuldung und ihre Grenzen}

Zu unterscheiden sind die jährliche \textbf{Neuverschuldung} (Kredite zum Ausgleich eines
Haushaltsdefizits) und die \textbf{Gesamtverschuldung} des Staates. Hohe Schulden und
Zinsverpflichtungen beschränken die Handlungsfähigkeit künftiger Politik. Juristische
Grenzen: die \textbf{Schuldenbremse} im Grundgesetz (Art.\,109 GG, seit 2009) und der
\emph{europäische Fiskalpakt} (2013), der die Mitgliedsländer verpflichtet, die
Schuldenbremse national zu regeln.

\section{Wachstumspolitik}

Wachstumspolitik umfasst alle staatlichen Maßnahmen für ein \emph{angemessenes und
stetiges} Wirtschaftswachstum. Zu unterscheiden sind \textbf{quantitatives Wachstum}
(reine Zunahme des realen BIP bzw. Pro-Kopf-Einkommens) und \textbf{qualitatives Wachstum}
(zusätzlich Schonung der Umwelt und Verbesserung der Lebensqualität; Messkonzepte: soziale
Indikatoren der OECD, Net Economic Welfare). Ansatzpunkte der Wachstumspolitik:

\begin{itemize}
  \item Förderung von \textbf{Investitionen} (besseres Investitionsklima, Gewinnerwartung),
  \item Beschleunigung des technischen Fortschritts durch Förderung von
    \textbf{Forschung und Entwicklung},
  \item Förderung von \textbf{Bildung}, um das Angebot an qualifizierten Arbeitskräften zu
    erhöhen,
  \item im weiteren Sinne auch die \textbf{Wettbewerbspolitik} (funktionierender
    Leistungswettbewerb) und die wirtschaftliche Infrastruktur.
\end{itemize}

\section{Tarifpolitik}

Die Tarifpolitik umfasst alle Fragen rund um Verhandlung und Abschluss von
\textbf{Tarifverträgen} zwischen den Tarifparteien --- in Deutschland in erster Linie
\emph{Arbeitgeberverbände} und \emph{Gewerkschaften}; auch ein einzelner Arbeitgeber kann
Tarifvertragspartei sein (\emph{Haustarifvertrag}).

\begin{merksatz}[title={\bfseries Tarifautonomie (Art. 9 Abs. 3 GG)}]
Die Tarifpartner sind von Weisungen des Staates \textbf{unabhängig}. Der Abschluss von
Tarifverträgen liegt ausschließlich in den Händen der Tarifparteien; der Staat ist zur
Neutralität verpflichtet. Staatliche Vorgaben sind nur als \emph{Gesetze} zulässig
(z.\,B. Bundesurlaubsgesetz, Entgeltfortzahlungsgesetz, Arbeitszeitgesetz,
Mindestlohngesetz) --- sie setzen Mindestansprüche, die nicht unterschritten werden dürfen.
\end{merksatz}

\subsection{Tarifvertragsarten}

\begin{itemize}
  \item \textbf{Lohn- und Gehalts(rahmen)tarifverträge:} regeln die Entgeltfindung ---
    Höhe des Entgelts, Eingruppierung, Zuschläge, Zulagen;
  \item \textbf{Manteltarifverträge:} regeln die arbeitsrechtlichen Rahmenbedingungen ---
    Verteilung der Wochenarbeitszeit, Pausen, Urlaub, Schichtarbeit, Freistellungen,
    Altersversorgung (meist mit langer Laufzeit).
\end{itemize}

\subsection{Lohnpolitik: Kaufkrafttheorie vs. produktivitätsorientierte Lohnpolitik}

\begin{itemize}
  \item Die \textbf{Gewerkschaften} argumentieren \emph{nachfrageorientiert}
    (\textbf{Kaufkrafttheorie}): Höhere Löhne steigern die Nachfrage, regen Produktion und
    Investitionen an und verbessern die Beschäftigung.
  \item Die \textbf{Arbeitgeber} argumentieren \emph{angebotsorientiert}: Löhne sind ein
    Kostenfaktor, der die (internationale) Wettbewerbsfähigkeit gefährdet. Sie fordern eine
    \textbf{produktivitätsorientierte Lohnpolitik}: Lohnsteigerungen nicht höher als der
    Produktivitätsfortschritt --- das verhindert Lohnkosteninflation und lohnkostenbedingte
    Arbeitslosigkeit.
\end{itemize}

Lohnerhöhungen wirken doppelt: Sie erhöhen die Produktionskosten
(\emph{Kostenerhöhungseffekt}), können aber durch höhere Produktivität der eingesetzten
Arbeit kompensiert werden (\emph{Kosteneinsparungseffekt}). Volkswirtschaftlich schädlich
sind überhöhte Abschlüsse dann, wenn Unternehmen mit Preiserhöhungen, Rationalisierung oder
Produktionsverlagerung ins Ausland reagieren.

\section{Arbeitsmarktpolitik}

Der Arbeitsmarkt ist kein „freier Markt`` im Sinne der Preisbildung: Wegen der hohen
sozialen Bedeutung des Lohns besteht Konsens, dass der Arbeitsmarkt nicht vollständig dem
freien Spiel der Marktkräfte überlassen wird. Zuständig für die Arbeitsmarktpolitik im
engeren Sinne ist v.\,a. die \textbf{Bundesagentur für Arbeit} (Nürnberg) mit ihren
Regionalagenturen, Arbeitsagenturen und Jobcentern.

\begin{itemize}
  \item \textbf{Passive Arbeitsmarktpolitik:} mildert die wirtschaftlichen Folgen
    eingetretener oder drohender Arbeitslosigkeit --- \emph{Lohnersatzleistungen} wie
    Arbeitslosengeld I und II (Bürgergeld), Insolvenzgeld, Kurzarbeitergeld.
  \item \textbf{Aktive Arbeitsmarktpolitik:} versetzt Erwerbspersonen in die Lage, im
    Arbeitsprozess zu bleiben oder schnell wieder einzutreten --- Beratung und
    \emph{Vermittlung}, Aktivierung und berufliche Eingliederung (Trainings- und
    Qualifizierungsmaßnahmen), Förderung der \emph{Aus- und Weiterbildung}, Umschulung und
    berufliche Rehabilitation, \emph{Existenzgründungsförderung}, Mobilitätshilfen (Reise-
    und Umzugskosten).
\end{itemize}

Leitgedanke ist heute \textbf{„Fördern und Fordern``}: Arbeitslose müssen sich aktiv
bemühen und zumutbare Jobs annehmen, sonst drohen Kürzungen der Lohnersatzleistungen.
Kritikpunkte an der Arbeitsmarktpolitik: geschönte Arbeitslosenquoten (Geförderte fallen
aus der Statistik), teils ineffiziente Qualifizierungsmaßnahmen, Verschlechterung der
Qualität der Arbeitsverhältnisse (Befristung, Leiharbeit, Scheinselbstständigkeit).

\section{Umweltpolitik}

\subsection{Prinzipien der Umweltpolitik}

\begin{itemize}
  \item \textbf{Verursacherprinzip:} Die Kosten für Vermeidung, Beseitigung oder Ausgleich
    von Umweltbelastungen trägt der Verursacher --- externe Kosten werden
    \emph{internalisiert} (gehen in Kostenrechnung und Preise ein); Beispiele:
    Rücknahmepflicht für Umverpackungen, Abwasserabgabe.
  \item \textbf{Gemeinlastprinzip:} Ist kein Verursacher feststellbar, übernimmt der Staat
    die Maßnahmen --- finanziert über Steuern (z.\,B. kommunale Kläranlagen, Mülldeponien);
    Nachteil: keine Anreize für umweltbewusstes Handeln.
  \item \textbf{Vorsorgeprinzip:} Umweltgefahren vorausschauend vermeiden --- „besser
    Vorsorge statt Nachsorge`` (viele Schäden, z.\,B. das Ozonloch, sind kaum reparabel).
  \item \textbf{Kooperationsprinzip:} Zusammenarbeit von Bürgern, Unternehmen und Staat.
  \item \textbf{Nachhaltigkeitsprinzip} (sustainable development, ursprünglich aus der
    Forstwirtschaft): künftige Generationen nicht schlechterstellen --- erneuerbare
    Ressourcen nur im Umfang ihrer Regenerationsfähigkeit nutzen, nicht erneuerbare sparsam
    nutzen, Emissionen an der Aufnahmefähigkeit der Ökosysteme ausrichten (vgl. Agenda~21).
\end{itemize}

\subsection{Instrumente der staatlichen Umweltpolitik}

\begin{itemize}
  \item \textbf{Nicht fiskalische Instrumente:} \emph{Aufklärung und Information};
    \emph{Auflagen} (Ordnungsrecht): umweltbezogene Ge- und Verbote, z.\,B. Grenzwerte für
    Schadstoffausstoß, FCKW-Verbot, Umweltplaketten/Umweltzonen.
  \item \textbf{Fiskalische Instrumente:}
    \emph{Umweltabgaben} --- Steuern (Ökosteuer), Sonderabgaben (Abwasserabgabe), Gebühren
    (Abfall-/Abwassergebühren) verteuern umweltschädliches Verhalten;
    \emph{Umweltzertifikate} --- Umweltbelastung nur gegen Berechtigung: Ein Zertifikat
    gestattet eine bestimmte Menge Schadstoffausstoß pro Jahr; die Emissionsrechte werden
    börsenmäßig gehandelt (EU-Emissionshandel für CO\textsubscript{2});
    \emph{Umweltsubventionen} --- Zuschüsse oder Steuererleichterungen für
    umweltschützende Maßnahmen (z.\,B. Wärmedämmung).
\end{itemize}

\begin{eselsbruecke}
Die drei großen Instrumentenfamilien: \textbf{A-A-Z} --- \textbf{A}uflagen (Ordnungsrecht:
verbieten), \textbf{A}bgaben (verteuern), \textbf{Z}ertifikate (begrenzen und handeln).
\end{eselsbruecke}

\section{Nachfrage- vs. angebotsorientierte Wirtschaftspolitik}

Nachfrage- und Angebotstheorie liefern unterschiedliche idealtypische Konzepte zur
Bewältigung von Wachstums- und Beschäftigungskrisen. Sie unterscheiden sich darin, welcher
\emph{Marktseite} die größere Bedeutung beigemessen wird und welche \emph{Rolle der Staat}
spielen soll.

\subsection{Nachfragetheorie nach Keynes (Fiskalismus)}

Die Nachfragetheorie geht auf \textbf{John Maynard Keynes} (1883--1946) zurück und sieht
Wachstumsschwächen und Arbeitslosigkeit in einer \emph{unzureichenden gesamtwirtschaftlichen
Nachfrage} begründet. Da der Markt allein das Gleichgewicht nicht herstellt, soll der Staat
eine \textbf{aktive Rolle} übernehmen und für zusätzliche Nachfrage sorgen
(\emph{Symptombekämpfung}, kurzfristig): Stärkung der Massenkaufkraft durch Lohnerhöhungen
und staatliche Zuschüsse, Erhöhung des Staatskonsums durch öffentliche Ausgabenprogramme,
Ausweitung des öffentlichen Sektors --- notfalls über Verschuldung (\emph{Deficit
Spending}). Geldpolitisch soll die Zentralbank die Finanzpolitik unterstützen (Liquidität
bereitstellen, Zinsen niedrig halten); Fiskal- und Geldpolitik wirken \emph{antizyklisch}.
Anhänger heißen \textbf{Fiskalisten}. Grundlage der Lohnargumentation ist die
\emph{Kaufkrafttheorie}.

\textbf{Schwächen:} zeitliche Verzögerungen (\emph{Timelags}) zwischen Entscheidung und
Wirkung, Verdrängung privater Investitionen durch staatliche (\emph{Crowding-out}),
steigende Staatsquote sowie steigende Abgaben-, Steuer- und Schuldenlast.

\subsection{Angebotstheorie (Neoklassik, Monetarismus)}

Die Angebotstheorie geht auf die klassische Nationalökonomie zurück
(\textbf{Neoklassizismus}; prominenter Vertreter: \textbf{Milton Friedman}, 1912--2006).
Beschäftigung und Wachstum hängen demnach von den Bedingungen auf der
\emph{Angebotsseite} ab; der Staat soll sich \textbf{passiv} verhalten und nur die
\emph{Rahmenbedingungen} verbessern (\emph{Ursachenbekämpfung}, mittel- bis langfristig):
Kostendämpfung zur Erhöhung der Unternehmensrentabilität, Steuersenkungen, Abbau von
Subventionen und Regulierungen (Deregulierung), Verringerung des effizienzschwachen
Staatskonsums, Investitionsförderung, moderate Lohnabschlüsse. Motor der Entwicklung sind
die \emph{Investitionen der Unternehmen}: Bessere Renditeerwartungen führen zu
Investitionen und neuen Arbeitsplätzen. Die Forderung nach \emph{stetigem
Geldmengenwachstum} basiert auf dem \textbf{Monetarismus}: Zwischen Geldmenge und BIP
besteht eine enge Beziehung.

\textbf{Schwächen:} einseitige Betonung der Kosten --- Löhne sind auch ein
\emph{Nachfragefaktor}; eine bessere Gewinnsituation erhöht zwar die
Investitions\emph{fähigkeit}, nicht automatisch die Investitions\emph{bereitschaft} (bei
geringer Absatzerwartung bleiben Erweiterungsinvestitionen aus).

\begin{eselsbruecke}
\textbf{Keynes = Nachfrage = Staat aktiv = kurzfristig Symptome behandeln.}
\textbf{Angebot = Rahmenbedingungen = Staat passiv = langfristig Ursachen beheben.}
Merkhilfe: „\textbf{K}eynes \textbf{k}auft`` (Staat schafft Nachfrage) ---
„\textbf{F}riedman \textbf{f}reit die Märkte`` (Deregulierung).
\end{eselsbruecke}

\section{Außenwirtschaft: Freihandel oder Protektionismus}

Unter \textbf{Außenwirtschaft} versteht man die Gesamtheit aller wirtschaftlichen
Aktivitäten zwischen Wirtschaftssubjekten im Inland und im Ausland (Handel mit Gütern und
Dienstleistungen, Kapital- und Devisenverkehr; ausgewiesen in der \emph{Zahlungsbilanz}).
Gründe für internationalen Handel: unterschiedliche klimatische Bedingungen, ungleich
verteilte Rohstoff- und Energievorräte, unterschiedliche Spezialisierungen und
Technologieschwerpunkte --- Grundgedanke der \textbf{internationalen Arbeitsteilung}.

\subsection{Freihandel}

\textbf{Freihandel} ist der völlig unbehinderte internationale Güteraustausch --- der Staat
verzichtet auf marktbeeinflussende Maßnahmen. Vorteile: große Güterauswahl in verschiedenen
Qualitäten und Preislagen für die Verbraucher, \emph{Spezialisierungsvorteile} und
Kostensenkungen in der Produktion, Absatzmärkte im In- und Ausland, Innovationsförderung
durch internationalen Wettbewerb. Gewinner sind allerdings v.\,a. die starken,
international wettbewerbsfähigen Unternehmen.

\subsection{Protektionismus: tarifäre und nichttarifäre Handelshemmnisse}

Greift der Staat zum Schutz der einheimischen Wirtschaft in den freien Handel ein, spricht
man von \textbf{Protektionismus} (Motive: Schutz vor Billigimporten, Vermeidung von
Monostrukturen und Auslandsabhängigkeit).

\begin{itemize}
  \item \textbf{Tarifäre Handelshemmnisse} (preispolitisch): \emph{Importzölle} (größte
    Bedeutung --- verbessern die Wettbewerbsposition der einheimischen Wirtschaft und
    bringen Staatseinnahmen), Steuertarife, \emph{Exportsubventionen} (inländische Produkte
    können im Ausland günstiger angeboten werden). Innerhalb der EU werden keine Zölle
    erhoben; gegenüber Drittländern gilt ein einheitliches EU-Vorgehen (keine autonome
    deutsche Zollpolitik mehr).
  \item \textbf{Nichttarifäre Handelshemmnisse} (mengenpolitisch/administrativ):
    \emph{Importkontingente} (mengenmäßige Beschränkungen), \emph{Ein- und Ausfuhrverbote}
    (schärfste Form, z.\,B. Rüstungsgüter, Drogen), \emph{Embargos} (außen- oder
    sicherheitspolitisch angeordnete Beschränkungen, i.\,d.\,R. Umsetzung internationaler
    Sanktionen von UNO oder EU), administrative Erschwernisse (Abfertigungs- und
    Genehmigungsverfahren), technische Normen, Gesundheits-, Sicherheits- und
    Verbraucherschutzstandards --- letztere mit stark zunehmender Bedeutung.
\end{itemize}

\begin{merksatz}
Rechtsgrundlagen in Deutschland: \textbf{Außenwirtschaftsgesetz (AWG)} und
Außenwirtschaftsverordnung (AWV); das \textbf{BAFA} (Bundesamt für Wirtschaft und
Ausfuhrkontrolle) informiert über Verbote, Genehmigungen und Beschränkungen.
\textbf{Globalisierung} = weltweiter Austausch von Gütern, Dienstleistungen, Informationen
und Kapital sowie weltweite Vernetzung von Unternehmen. Die EU ist Mitglied der
\textbf{WTO} und schließt Freihandelsabkommen (z.\,B. CETA mit Kanada, JEFTA mit Japan).
\end{merksatz}

\section{Besonderheiten der EU}

Die Integration begann 1957 mit der \emph{Europäischen Wirtschaftsgemeinschaft (EWG)}
(Römische Verträge, sechs Gründungsstaaten); mit dem Vertrag von \textbf{Maastricht} (1993)
entstand die \emph{Europäische Union (EU)}.

\subsection{Der Europäische Binnenmarkt: die vier Freiheiten}

Ein \textbf{Binnenmarkt} ist ein Wirtschaftsraum ohne innere Grenzen (gegründet 1993). Er
ist durch \textbf{vier Freiheiten} gekennzeichnet:

\begin{center}
\begin{tabular}{p{4.2cm} p{9.6cm}}
\toprule
\textbf{Freiheit} & \textbf{Kennzeichen}\\
\midrule
Freier Personenverkehr & Abbau der Grenzkontrollen (Schengen), Aufenthalts- und Niederlassungsfreiheit, freie Arbeitsplatzwahl, Anerkennung von Berufsabschlüssen\\
Freier Warenverkehr & keine Zölle oder mengenmäßigen Beschränkungen, Harmonisierung von Umsatz-/Verbrauchssteuern und technischen Normen\\
Freier Dienstleistungsverkehr & Selbstständige und freie Berufe dürfen EU-weit anbieten und sich niederlassen; Liberalisierung bei Banken und Versicherungen\\
Freier Geld- und Kapitalverkehr & Öffnung der Finanzmärkte, erleichterter Zahlungsverkehr (SEPA), Aufhebung von Kapitalverkehrsbeschränkungen\\
\bottomrule
\end{tabular}
\end{center}

\textbf{Chancen:} größerer Absatzmarkt, sinkende Stückkosten durch höhere Stückzahlen und
Rationalisierung, vielfältigeres und günstigeres Güterangebot, stabilisierende Wirkung auf
das Preisniveau. \textbf{Risiken/Nachteile:} Verlagerung von Produktionsstätten ins
kostengünstigere EU-Ausland (Arbeitsplatzverluste), Bürokratiekosten, aus Sicht der
Kritiker „Entmachtung`` der Einzelstaaten.

\begin{eselsbruecke}
Die vier Freiheiten: \textbf{P-W-D-K} --- \textbf{P}ersonen, \textbf{W}aren,
\textbf{D}ienstleistungen, \textbf{K}apital. Merkspruch: „\emph{Praktisch wandern durch
Kontinente}``.
\end{eselsbruecke}

\subsection{Europäische Wirtschafts- und Währungsunion (EWWU)}

Die Europäische Währungsunion ist Teil der EWWU; die Teilnehmerstaaten bilden die
\textbf{Eurozone}. Nach dem Maastrichter Vertrag dürfen nur EU-Staaten den Euro einführen,
die ihre Stabilität durch die \textbf{Konvergenzkriterien} bewiesen haben:

\begin{enumerate}
  \item \textbf{Inflationsrate:} im Jahr vor dem Beitritt höchstens 1,5~Prozentpunkte über
    dem Mittelwert der drei preisstabilsten EU-Staaten;
  \item \textbf{Nettoneuverschuldung:} nicht dauerhaft mehr als \textbf{3~\%} des BIP;
  \item \textbf{Gesamtschuldenstand:} i.\,d.\,R. nicht mehr als \textbf{60~\%} des BIP;
  \item \textbf{Zinssatz für langfristige Kredite:} höchstens 2~Prozentpunkte über dem
    Mittelwert der drei preisstabilsten EU-Staaten;
  \item \textbf{Wechselkursstabilität:} zwei Jahre stabiler Außenwert ohne Abwertung
    (Teilnahme am Wechselkursmechanismus WKM~II).
\end{enumerate}

Die Einhaltung wurde in der Vergangenheit nicht konsequent überprüft; mit der Euro- und
Staatsschuldenkrise reformierten EU-Parlament und Finanzminister 2011 den europäischen
\emph{Stabilitäts- und Wachstumspakt} (härteres Vorgehen gegen „Haushaltssünder``; die
Referenzwerte 3~\% und 60~\% blieben erhalten).

\subsection{Vor- und Nachteile der gemeinsamen Währung}

\begin{itemize}
  \item \textbf{Vorteile:} keine Wechselkursrisiken und Umtauschkosten im Euroraum,
    Preistransparenz (direkter Preisvergleich), Förderung von Handel, Investitionen und
    Reiseverkehr, starke Währung im Weltmaßstab.
  \item \textbf{Nachteile:} Verzicht auf eigenständige nationale Geld- und
    Wechselkurspolitik (keine Abwertung zur Stärkung der Wettbewerbsfähigkeit möglich),
    einheitlicher Leitzins für unterschiedlich entwickelte Volkswirtschaften,
    Haftungs- und Transferrisiken bei Schuldenkrisen einzelner Mitglieder.
\end{itemize}

\section{Formeln \& Rechenbeispiele}

Die folgenden drei Berechnungen sind Prüfungsklassiker --- alle folgen demselben Muster:
\emph{Teil durch Ganzes mal 100}. Prägen Sie sich Formel und Rechenweg anhand der
einfachen Zahlenbeispiele ein.

\begin{merksatz}[title={\bfseries Formel 1: Wirtschaftswachstum (Veränderungsrate des realen BIP)}]
\[
\text{Wachstumsrate in \%} =
\frac{\text{BIP}_{\text{neu}} - \text{BIP}_{\text{alt}}}{\text{BIP}_{\text{alt}}} \times 100
\]
Immer mit dem \emph{realen} BIP rechnen (Preise eines Basisjahres) --- sonst misst man
Inflation statt echtem Wachstum.
\end{merksatz}

\begin{beispiel}
Das reale BIP steigt von 3.000~Mrd.~Euro (Vorjahr) auf 3.060~Mrd.~Euro.\\[0.3em]
\textbf{Schritt 1:} Veränderung berechnen: $3.060 - 3.000 = 60$~Mrd.~Euro.\\
\textbf{Schritt 2:} Durch den \emph{alten} Wert teilen: $60 \div 3.000 = 0{,}02$.\\
\textbf{Schritt 3:} Mal 100: $0{,}02 \times 100 = \textbf{2{,}0~\%}$ Wirtschaftswachstum.
\end{beispiel}

\begin{merksatz}[title={\bfseries Formel 2: Inflationsrate (Veränderungsrate des Preisindex)}]
\[
\text{Inflationsrate in \%} =
\frac{\text{Preisindex}_{\text{neu}} - \text{Preisindex}_{\text{alt}}}{\text{Preisindex}_{\text{alt}}} \times 100
\]
Grundlage ist der Verbraucherpreisindex (in der Eurozone: HVPI), der die Preisentwicklung
eines repräsentativen Warenkorbs misst.
\end{merksatz}

\begin{beispiel}
Der Verbraucherpreisindex steigt von 110 auf 115,5 Punkte.\\[0.3em]
\textbf{Schritt 1:} Veränderung: $115{,}5 - 110 = 5{,}5$ Punkte.\\
\textbf{Schritt 2:} Durch den \emph{alten} Indexstand teilen: $5{,}5 \div 110 = 0{,}05$.\\
\textbf{Schritt 3:} Mal 100: $0{,}05 \times 100 = \textbf{5{,}0~\%}$ Inflationsrate.\\[0.3em]
\emph{Typischer Fehler:} durch 100 statt durch den Basiswert 110 teilen --- das ergäbe
fälschlich 5,5~\%.
\end{beispiel}

\begin{merksatz}[title={\bfseries Formel 3: Arbeitslosenquote}]
\[
\text{Arbeitslosenquote in \%} =
\frac{\text{registrierte Arbeitslose}}{\text{alle zivilen Erwerbspersonen}} \times 100
\]
Zivile Erwerbspersonen = Erwerbstätige \emph{plus} registrierte Arbeitslose (also alle,
die arbeiten oder arbeiten wollen).
\end{merksatz}

\begin{beispiel}
In einem Land sind 2,3~Mio. Personen arbeitslos gemeldet; es gibt 43,7~Mio. Erwerbstätige.\\[0.3em]
\textbf{Schritt 1:} Erwerbspersonen bestimmen: $43{,}7 + 2{,}3 = 46$~Mio.\\
\textbf{Schritt 2:} Arbeitslose durch Erwerbspersonen teilen: $2{,}3 \div 46 = 0{,}05$.\\
\textbf{Schritt 3:} Mal 100: $0{,}05 \times 100 = \textbf{5{,}0~\%}$ Arbeitslosenquote.
\end{beispiel}

\begin{merksatz}[title={\bfseries Anwendung: Konvergenzkriterien nachrechnen (3~\% / 60~\%)}]
\[
\text{Neuverschuldungsquote} = \frac{\text{jährliche Neuverschuldung}}{\text{BIP}} \times 100
\qquad
\text{Schuldenstandsquote} = \frac{\text{Gesamtschulden}}{\text{BIP}} \times 100
\]
Grenzwerte: Neuverschuldung max. \textbf{3~\%}, Gesamtschuldenstand max. \textbf{60~\%}
des BIP.
\end{merksatz}

\begin{beispiel}
Ein Beitrittskandidat hat ein BIP von 500~Mrd.~Euro, eine Neuverschuldung von
20~Mrd.~Euro und Gesamtschulden von 280~Mrd.~Euro.\\[0.3em]
\textbf{Schritt 1:} Neuverschuldung: $20 \div 500 \times 100 = 4{,}0~\%$ ---
\emph{über} der 3-\%-Grenze~$\Rightarrow$ Kriterium \textbf{verletzt}.\\
\textbf{Schritt 2:} Schuldenstand: $280 \div 500 \times 100 = 56{,}0~\%$ ---
\emph{unter} der 60-\%-Grenze~$\Rightarrow$ Kriterium \textbf{erfüllt}.\\
\textbf{Ergebnis:} Das Land erfüllt die fiskalischen Konvergenzkriterien insgesamt
\emph{nicht} (beide Grenzwerte müssen eingehalten werden).
\end{beispiel}

\begin{merksatz}[title={\bfseries Wichtig für die Prüfungsvorbereitung (Block 1.1 komplett)}]
Sie sollten sicher können: den \emph{Konjunkturzyklus} mit Indikatoren, das Ziel und die
\emph{drei Instrumente der EZB} samt \emph{Leitzinsen} und der Wirkungskette einer
Leitzinsänderung, die \emph{antizyklische Fiskalpolitik} (Boom vs. Rezession) mit
automatischen Stabilisatoren, die Gegenüberstellung \emph{nachfrage- vs.
angebotsorientiert}, \emph{tarifäre vs. nichttarifäre Handelshemmnisse}, die \emph{vier
Freiheiten} des Binnenmarkts, die \emph{Konvergenzkriterien} (3~\%~/ 60~\%) sowie die
\emph{Formeln} für Wachstumsrate, Inflationsrate und Arbeitslosenquote (Abschnitt~11).
Wiederholen Sie außerdem
Lektion~1 (Preisbildung, Marktformen) und Lektion~2 (Wettbewerb, Staatseingriffe, VGR,
magisches Viereck/Sechseck) --- das Quiz dieser Lektion und das Prüfungstraining helfen dabei.
\end{merksatz}

\vspace{1em}
\hrule
\vspace{0.8em}
\noindent\textsf{\small Quellen: Rahmenplan Wirtschaftsbezogene Qualifikationen (DIHK), Abschnitte 1.1.3.2 und 1.1.4;
Textband Volkswirtschaft (DIHK-Bildungs-GmbH), Kap.~3.2 und Kap.~4, S.~39--64. Nächster Termin (Sa, 05.09.2026):
Start des BWL-Teils --- Betriebliche Funktionen~I.}

\end{document}
